% Training and evaluation procedure for the four response-delay models.
% Terminology follows Agarwal et al., "Deep Reinforcement Learning at the Edge of the
% Statistical Precipice" (NeurIPS 2021) and the rliable library. US spelling to match.

\documentclass[11pt,a4paper]{article}

\usepackage[margin=2.5cm]{geometry}
\usepackage{amsmath}
\usepackage{booktabs}
\usepackage[hidelinks]{hyperref}

\newcommand{\iqm}{\textsc{iqm}}
\newcommand{\nocr}{\textsc{no-cr}}

\title{Training and Evaluation Procedure}
\author{}
\date{}

\begin{document}
\maketitle

%% ---- body starts here ----

\section*{Objective}

Each of the four policies studied in this work is trained under a different
model of how long pilots take to act on an advisory. All four are then evaluated
in simulated environments in which pilots respond after a lognormally
distributed delay, and the mean of that delay is varied across a range of
values. Because the same distribution is used at test time for every policy, any
difference between them can be attributed to the delay model each was trained
under rather than to the conditions under which they were compared.

The work proceeds in three steps. Step 1 trains the policies. Step 2 evaluates
each trained policy over a fixed set of episodes and records the results in
tables. Step 3 analyzes those tables with the tools of \texttt{rliable}.

\section{Step 1 --- Training}

Four training conditions are used, listed in Table~\ref{tab:models}. Each
condition is trained five times with a different random seed, which gives twenty
trained policies in total.

\begin{table}[h]
\centering
\begin{tabular}{@{}lll@{}}
\toprule
Label & Delay model & Behavior \\
\midrule
N & None          & Advisories take effect immediately. \\
D & Deterministic & Every advisory takes effect exactly 30 seconds after it is issued. \\
G & Geometric     & The pilot has the same chance of acting in every second, so the \\
  &               & delay is memoryless and its tail is unbounded. \\
L & Lognormal     & Delays cluster around a typical value with a longer tail toward \\
  &               & slow responses. \\
\bottomrule
\end{tabular}
\caption{The four training conditions. The three delayed conditions share a mean
response time of 30 seconds and differ only in the shape of the distribution
around that mean.}
\label{tab:models}
\end{table}

Training more than one policy per condition is what makes the statistical
analysis in Step 3 possible. A single training run yields one value for each
measurement, and that value gives no indication of how much of the result
reflects the delay model and how much reflects the particular seed the run
happened to use. Reinforcement learning is known to vary considerably between
seeds, so a difference between two conditions is only meaningful if it is larger
than the variation between runs of the same condition. The five runs per
condition provide the estimate of that variation.

The same five seeds are used in all four conditions. Pairing the conditions in
this way means that corresponding runs begin from the same network
initialization and see the same sequence of training scenarios, so that the
delay model is the only thing that differs between them. The seed values
themselves are spaced widely apart, because the parallel workers within a run
are seeded by increments from the run seed. Closely spaced run seeds would
otherwise cause two runs to share part of their training scenarios, and the
runs would no longer be independent.

\begin{table}[h]
\centering
\begin{tabular}{@{}ll@{}}
\toprule
Runs per condition & Five, or ten if the compute budget allows \\
Training length    & Ten million environment steps per run \\
Seeds              & The same five, widely spaced, are used in every condition \\
Scenarios          & Drawn from a training pool that is disjoint from the test pool \\
Policy reported    & The policy that training ended on \\
\bottomrule
\end{tabular}
\caption{Training protocol. The total cost equals that of training four single
policies for fifty million steps each, so the budget is redistributed rather
than increased.}
\label{tab:training}
\end{table}

Before committing to all twenty runs, one pilot run per condition should be
trained and its learning curve inspected. If performance has clearly leveled off
before the budget is exhausted, the budget is adequate. If it has not, the
training length should be increased while keeping five seeds per condition. The
number of seeds should not be reduced to pay for longer runs.

\paragraph{Output of Step 1.} Twenty trained policies, one for each combination
of the four conditions and five seeds.

\section{Step 2 --- Validation}

Each of the twenty trained policies is evaluated in environments where the pilot
response delay follows a lognormal distribution. Seven mean delays are used,
ranging from zero to two minutes, with zero corresponding to an environment
with no delay at all.

Every evaluation uses the same fixed set of one hundred scenarios, drawn from a
test pool that does not overlap with the scenarios seen during training. Holding
the scenarios constant across all conditions, seeds and delay levels ensures
that differences in the results come from the policy or from the delay, and not
from one policy happening to be given an easier set of sectors than another.

One reference policy is evaluated alongside the twenty trained ones: the \nocr{}
policy, which never issues an advisory. Since it transmits nothing, the pilot
response delay has no effect on it, and it needs to be evaluated only once. It
appears in the results as the level of performance obtained without any conflict
resolution at all.

The complete grid comprises four conditions, seven delay levels and five seeds,
which is 140 evaluation runs. Together with the reference policy this amounts to
roughly fourteen thousand episodes.

\begin{table}[h]
\centering
\begin{tabular}{@{}ll@{}}
\toprule
Indicator & Preferred direction \\
\midrule
Losses of separation per flight hour   & lower \\
Fraction of time in loss of separation & lower \\
Arrival rate, as on-route exits        & higher \\
Exit deviation from the planned route  & lower \\
Heading changes per flight hour        & lower \\
Speed changes per flight hour          & lower \\
Total advisories issued                & lower \\
Episode reward                         & higher \\
\bottomrule
\end{tabular}
\caption{The eight performance indicators recorded for every episode. All are
kept in their original physical units and are never rescaled.}
\label{tab:kpis}
\end{table}

\paragraph{Output of Step 2.} One table per evaluation run, containing one row
per episode. Each row records the training condition, the training seed, the
mean test delay and the scenario identifier, together with the eight performance
indicators listed in Table~\ref{tab:kpis}. A small number of diagnostic
quantities, such as the response delay actually observed, are also retained as a
check that the delay model behaved as configured. No conclusions are drawn from
them.

\section{Step 3 --- Analysis}

The tables produced in Step 2 are combined so that, for each combination of
training condition, test delay and performance indicator, the measurements form
a matrix with one row for each training run and one column for each scenario.
This is the layout expected by the statistical tools described below, in which
the training runs play the role of independent repetitions of the experiment and
the scenarios play the role of separate tasks.

Table~\ref{tab:matrix} gives an example. It holds the losses of separation per
flight hour recorded by condition L when flown at a mean test delay of 45
seconds. Each row is one of the five training runs and each column is one of the
hundred evaluation scenarios, so the matrix contains five hundred measurements.
One matrix of this kind exists for every combination of the four conditions,
seven delay levels and eight indicators, which is 224 in total.

\begin{table}[h]
\centering
\begin{tabular}{@{}lccccc@{}}
\toprule
      & Scenario 1 & Scenario 2 & Scenario 3 & $\cdots$ & Scenario 100 \\
\midrule
Run 1 & 0.00 & 0.14 & 0.07 & $\cdots$ & 0.05 \\
Run 2 & 0.06 & 0.11 & 0.09 & $\cdots$ & 0.00 \\
Run 3 & 0.00 & 0.19 & 0.06 & $\cdots$ & 0.08 \\
Run 4 & 0.05 & 0.12 & 0.11 & $\cdots$ & 0.04 \\
Run 5 & 0.00 & 0.16 & 0.08 & $\cdots$ & 0.06 \\
\bottomrule
\end{tabular}
\caption{An example score matrix, holding losses of separation per flight hour
for one condition at one test delay. Scenario 2 is harder than scenario 1 for
every run, which is why the scenarios are treated as strata: resampling within a
column preserves those differences in difficulty instead of averaging them
away.}
\label{tab:matrix}
\end{table}

A matrix of this form is the only thing the analysis is given. The tools applied
to it are those of \texttt{rliable}, the library accompanying Agarwal et al.;
the definitions used here are implemented directly, so no external dependency is
needed. Each tool returns a summary together with its uncertainty.

\begin{table}[h]
\centering
\begin{tabular}{@{}lll@{}}
\toprule
Tool & Given & Returns \\
\midrule
Interquartile mean         & one matrix  & a number, with an interval \\
Performance profile        & one matrix  & a curve over thresholds, with a band \\
Probability of improvement & two matrices & a number, with an interval \\
\bottomrule
\end{tabular}
\caption{What each tool is handed and what it gives back.}
\label{tab:tools}
\end{table}

The uncertainty always comes from the same operation. The matrix is resampled by
drawing the five runs with replacement inside each column, leaving the set of
columns unchanged; the summary is recomputed on the resampled matrix; and this
is repeated fifty thousand times. The reported interval is the range covering
the middle 95 percent of the recomputed values.

\paragraph{Point estimate.} Each matrix is summarized by its interquartile mean,
which is the mean of the middle half of the measurements once the runs and
scenarios have been pooled. It is preferred to the ordinary mean because it is
not distorted by a small number of extreme episodes, and preferred to the median
because it uses more of the available data and therefore yields narrower
confidence intervals. The mean and the median are tabulated alongside it so that
a reader can confirm that the conclusions do not depend on which summary is
used.

\paragraph{Uncertainty.} Every reported value is accompanied by a 95 percent
confidence interval obtained by a stratified bootstrap. The procedure resamples
the training runs with replacement, independently within each scenario, so that
the scenarios act as strata and their relative difficulty is held fixed. Fifty
thousand resamples are used, with a fixed random seed so that the figures can be
reproduced exactly.

It matters that the interval is computed across training runs rather than across
the episodes of a single policy. An interval of the latter kind describes only
how much the scenario draw varies, and says nothing about how much the outcome
of training itself varies, which is the quantity the comparison between
conditions depends on.

\paragraph{Distributions.} A performance profile is plotted for episode reward.
For every threshold, it shows the fraction of runs across all scenarios that
achieved a reward above that threshold, together with bootstrap confidence
bands. The profile therefore presents the whole distribution of outcomes rather
than a single summary of it, which makes it possible to distinguish a policy
that is better in every scenario from one that is better only on average.

\paragraph{Comparisons.} The probability of improvement is reported for each
delayed condition against the undelayed baseline, and for the two stochastic
conditions against the deterministic one. It states how often one policy
outperforms another, independently of the size of the difference, and is the
appropriate measure when the question is whether training with a delay model
helps at all rather than by how much.

\paragraph{Output of Step 3.} Four figures.

\begin{enumerate}
  \item \emph{Degradation curves.} Eight panels, one for each indicator in
    Table~\ref{tab:kpis}, showing the interquartile mean in its original units
    against the mean test delay, with one curve per training condition,
    stratified bootstrap confidence bands, and the \nocr{} reference as a dashed
    horizontal line.
  \item \emph{Performance profiles} of episode reward, at a moderate and a
    severe mean test delay.
  \item \emph{Probability of improvement} of each condition over the baseline,
    plotted against the mean test delay.
  \item \emph{Learning curves} showing the interquartile mean of the training
    return across the five runs of each condition.
\end{enumerate}

Every caption states the number of training runs and the number of evaluation
scenarios behind the figure, and names the summary statistic that is plotted.

%% ---- body ends here ----

\end{document}
